\chapter{数值计算基础}
\label{chap:numerical-basics}

\section{引言}

\subsection{我们为什么需要数值计算？}

对于大多数高校的物理学专业，计算物理学通常都会作为一门必修课开设给那些完成了普通物理知识和基本编程语言学习的同学们。拿起本书的读者或许会产生这样的疑问：我们为什么需要数值计算？

同学们在学习了几年的物理学专业课程之后，提起物理学家，在脑海中往往会浮现这样的形象：他们面对任何问题，都能够迅速抓住主要因素，抽象出一个简单明了的物理模型并给出解析解，从而对现实世界进行预言或给出行为建议。抑或是在实验室中运用各种尖端手段，创造最纯净、最理想的物理条件，检验那些描述基本粒子运动的最简洁优美的方程。至于说在考虑了现实中诸多繁杂的“次要”因素后能否做出，以及如何做出精确预言，从基本的万物理论中如何演生出生命自身以及我们周围的宏观世界，那貌似只是技术和工程问题，或者属于化学家、生物学家和材料学家们所关心的范畴。物理学家所面对的永远都是最干净的研究对象，一支笔加一沓纸就能轻松应付。基于这个想法，编程和数值计算便和立志于从事物理学研究的同学们关系不大，似乎不应当成为关注的重点。

倘若你对上一段话表示认同，那么恭喜你，你所接受的基础物理教育是很成功的。物理学教会我们透过现象看本质，让我们笃信我们周围的世界是可以被观测和理解的。自 Galileo 和 Newton 以来，物理学在这样最理想、最经典的研究范式下取得了长足的进展。小到夸克、胶子，大到星系、宇宙，前辈们掌握了不同尺度上的物质世界的特性及其演化规律。

然而，这种研究范式并不总是有效。可以解析求解的物理系统其实是非常稀少的，以至于有戏言称理论物理学家只会处理平面波和谐振子。譬如，三个仅通过引力相互作用的质点所构成的系统，即三体问题，自百余年前被提出以来，人们至今仍不敢说完全掌握了其特性，而需要数值工具的辅助。近年来，对此系统的探讨依旧还在取得有深刻意义的研究成果。\footnote{Ginat Y B and Perets H B. Analytical, statistical approximate solution of dissipative and nondissipative binary-single stellar encounters. \textit{Phys. Rev. X}, 2021, 11: 031020.}

在本章末的大作业中，读者将会尝试通过数值实验去探讨一个更简单的系统：迭代式
\[
x_{n+1}=r x_n(1-x_n),
\]
它描述了很多不同现象背后的非线性动力学行为。

当所参与的粒子数和相互作用进一步增多，以至于不同物质层级之间的界限变得模糊时，问题将变得愈发复杂：如何用构成中子和质子的夸克之间的强相互作用力，来解释中子和质子之间的核力？如何用构成原子、分子的电子与原子核之间的电磁力，来解释原子、分子之间的共价键、离子键、氢键和 van der Waals 力？我们的物质世界基于包含概率和塌缩的量子力学，为何日常生活表现得如同由确定性的经典力学来描述？微观粒子的运动方程满足时间反演对称性，在宏观层面上是怎样出现单向的时间箭头和熵增的？生命系统又是如何克服熵增，不断有序地演化发展的？理论或许可以提供一些洞见，但如果能够通过数值方法给出系统在不同层级之间是如何过渡的，人们对这些问题的理解自然会更加深刻。

除了在理论上求解难以解析处理的物理系统，如今数值方法对于实验的推动也是不可或缺的。我们早已过了在坐标纸上记录实验数据拟合直线的时代，诸如高能物理实验和天文观测所积累的海量信号都对数值工具提出了非常高的要求。从中找寻新物理的蛛丝马迹，需要非常扎实的概率统计和数值分析的基础。不过，由于编者的背景所限，本书将不会对这部分内容过多涉及。

事实上，如图\ref{fig:three-tools}所示，数值计算与理论分析、科学实验已构成现代科学与技术研究的三大手段，数值计算的重要应用早已从科学研究延伸到工程应用、经济学、人文和社会科学等领域，全面推动着科技进步和人类文明的发展。具体来说，对于科学与工程应用领域，数值计算涉及读者所知晓的众多学科，包括物理学、化学、材料科学、生命科学、医学、计算机图形学、力学、天文学、大气和海洋科学、核科学及应用、航天航海、环境科学、工程设计与应用等。数值计算在很多实际问题中起着不可替代的作用，例如核武器的理论设计、气象预报与灾害预警、航天飞行器的轨道预设、汽车安全性的碰撞试验、新材料和新药物的设计等。总之，大规模的科学计算涉及科学研究、国民经济、国防建设和社会发展的方方面面，与其相关的软件和硬件发展水平是一个国家综合国力的重要标志之一。例如，“TOP500”榜单\footnote{\url{https://www.top500.org/}}是对全球已建成的超级计算机运算速度（算力）的权威排行榜，始于1993年，由国际组织“TOP500”编制，每半年发布一次。该榜单以实测 LINPACK benchmark 性能基准进行排名。从2016年开始，我国超级计算机进入“TOP500”榜单的数量，除了2017年6月稍下滑位居第二外，一直稳定在世界第一的位置（近年来，中国有很多超级计算机都不再参与这个排行榜了）。但是，我们必须清醒地认识到我国在硬件和软件自主化的道路上仍然任重道远。

\begin{figure}[htbp]
\centering
\begin{tikzpicture}[>=Stealth, node distance=18mm and 22mm,every node/.style={font=\small}]
\node[draw,ellipse,minimum width=30mm,minimum height=11mm] (num) {数值计算};
\node[draw,ellipse,below left=of num,minimum width=34mm,minimum height=11mm] (theory) {理论分析};
\node[draw,ellipse,below right=of num,minimum width=34mm,minimum height=11mm] (exp) {科学实验};
\draw[<->,thick] (num)--(theory);
\draw[<->,thick] (num)--(exp);
\draw[<->,thick] (theory)--(exp);
\end{tikzpicture}
\caption{现代科学、工程与技术研究相辅相成的三大手段}
\label{fig:three-tools}
\end{figure}

\subsection{计算物理学的特点与本书介绍}

本书讨论数值计算在物理学中的应用。计算物理学已成为一个重要的物理学分支，是利用计算机来解决物理问题或者分析物理实验结果的学科，有其相对的独立性，同时又是物理学、数学、计算机科学三者的交叉学科。本书假定读者掌握了微积分、线性代数，以及任何一门现代高级编程语言，对普通物理和四大力学有一些基本的了解。我们将试图以物理问题为主线，将数值计算的知识和技巧，融于对这些物理问题的算法设计和讨论之中。随着本书内容的推进，每一节中都会提供一种新的算法或技能，但读者往往需要将新习得的算法与之前掌握的某些其他算法有机结合在一起，才能解决所面对的新问题。另外，尽管本书主要以物理问题为主线，我们仍然尽力按照一定的逻辑关系来展现各种数值方法和算法之间的内在联系。

运用数值方法解决实际物理问题的过程与通常学过的方法并没有太大的差别，一般大致分成这样几个阶段：由于真实系统的复杂性，首先我们需要抓住问题的主要因素，忽略次要因素，从实际问题中抽象出物理模型，并归结为一个可求解的数学问题，也就是所谓的数学建模。然后，面对适定的数学问题，经过对特定问题的定性分析与考察，我们需要构造一套行之有效的数值计算方法。继而，我们要选择合适的编程语言，对所设计的数值算法实施程序设计。接下来，我们要全面调试整套程序，确保程序能够给出正确可靠的数值结果。最后，我们将所开发的整套程序，应用到当前的实际场景，运行程序得到问题的数值解，并谨慎地对结果进行误差分析和相关讨论，进而讨论数值结果所揭示的物理规律。

不同于计算机科学中常见的离散数学，物理学更关心那些可以连续变化的变量。因此，实数、解析函数、微分方程等都是计算物理问题中的常客。我们需要根据相应问题的特性设计算法、实现算法，并分析算法及其结果。所谓算法，不只是单纯的数学公式，而是一整套利用计算机来求解问题的方案和步骤，所有步骤都要归结为计算机可以执行的底层操作。实际上，计算机能够执行的只有有限字长的二进制逻辑运算，只不过通过硬件设计和底层编码构造了浮点数，实现了对实数及其基本四则运算的近似表示。高级编程语言中可以调用的诸多函数，例如三角函数、指数函数、对数函数等，都是通过内置的算法调用前述基本运算来近似实现的。

由于计算机的各种局限性，在最初原理和模型近似的基础之上，把一个实际问题抽象为适定的、可求解的数值问题时，也不得不采用一定的数学近似。我们使用有限长度的浮点数近似替代实数，也使用有限的四则运算过程近似实现复杂函数计算；更进一步说，当我们分析函数的行为，研究其导数和积分时，在计算机上不可能真正实现其定义所需要的极限过程，而会使用差分和有限求和来近似；在研究函数的微分方程或者积分方程问题时，也要考虑如何将无限维的函数空间近似成有限空间，从而转化为一个有限维的线性代数问题；求解偏微分方程需要完全掌握整个区域的初始条件和边界条件，而这在现实情况中几乎无法做到，我们可以将它们近似成前几阶“矩”相同而便于后续处理的函数形式。\footnote{在弹性力学领域，人们将这个近似的合理性建立在被称作 Saint-Venant 原理的基本假设上。Saint-Venant 原理虽然已经有大量实例验证，但至今还没有严格的数学证明。这一原理由法国科学家 Saint-Venant 在1855年提出，其基本内容是：分布于弹性体上一小块面积（或体积）内的荷载所引起的物体中的应力，在离荷载作用区域稍远的地方，基本上只同荷载的合力和合力矩有关，而荷载的具体分布只影响荷载作用区域附近的应力分布。}\jiaokan{此处原文作“1885年”。Saint-Venant 对这一原理的原始表述发表于1855年；可参见其同年出版的 \textit{De la torsion des prismes, avec des considérations sur leur flexion}。}我们通常要将无限量化作有限量，然后将过程中产生的误差减小到足够低，这个思想遍布于用计算机求解物理问题的每一个环节。

在本书中，我们会先在本章介绍数值计算的基本概念，并给出一些简单的例子让读者获得一些直观的体会。第二章会以质点的一维运动为引子，介绍对函数进行积分、求根、求极值以及求解微分方程的最基本的算法，让读者初步具有自行设计算法解决物理问题的能力。第三章将通过对线性电路网络和简谐振动的讨论，引入求解有限维矩阵问题的诸多数值算法，并为后续章节处理无穷维问题打下基础。第四章将前两章的内容做综合，讨论如何处理多自由度的非线性问题。第五章是我们第一次开始研究场的问题，主要思想便是如何将无穷维的函数空间离散化，转化成第三章中介绍过的数值方法能够处理的问题。第六章将处理更复杂的场，即量子场，探讨其相比经典场方程会多出哪些物理特性，以及如何进行数值处理。第七章的内容与前述章节大不相同，我们会处理基于概率的统计物理问题，并且探讨如何通过概率的手段来解决之前出现过的确定性问题。第八章将介绍近年来快速发展的机器学习算法，并讨论如何将其应用到物理学，解决实际的物理问题。

如前所述，我们在数值求解具体问题的过程中，总会引入各种各样的近似，而只要做近似，就会有误差：忽略实际问题中的次要因素建立数学模型时会产生模型误差，模型中输入的参数与实际物理量之间存在观测误差，浮点数做运算时会产生舍入误差，在数值算法中将无穷项截断至有限项会产生截断误差（这类误差是数值计算方法所固有的，因此又称为方法误差）。这些来自不同近似过程的误差在数值计算的进行中可能会被放大或压低，有可能某个一开始看上去微不足道的误差会造成最终结果的不可信，也有可能看起来完全是“错误”的中间过程会导向一个高精度的最终结果。一个理想的算法应当能保证所有来源所产生的误差对最终结果的影响在数量级上大致相同，从而在保证结果精度的同时不浪费太多的脑力、劳力和算力。当然这并不简单，实践中通常会重点考虑那几个最大的误差来源以求尽量将其压低。对误差的分析和控制将会贯穿本书的始终，也会贯穿在任何计算物理实践的全过程。

\subsection{数值计算的基本策略}

在着手设计算法来解决即使是最简单的物理问题之前，我们都必须对数值计算的基本特点有相当的了解，并牢记于心，在后续算法的设计和程序编写过程中，遵循一些基本的习惯和策略，才能使得数值求解该物理问题是有效的，得到的数值结果是可靠的。为此，我们以几个大家熟知的简单数学问题为例，来具体展示一下在实际编程过程中一个好的数值算法所需要考虑和注意的几个方面。

\subsubsection{等价与不等价}

同一个表达式，数学上往往有着多种等价的表示。比如说，一元二次方程
\begin{equation}
ax^2+bx+c=0,\qquad (a\ne 0),
\tag{1.1.1}
\end{equation}
其两个根可以写成
\begin{equation}
x_{\pm}=\frac{-b\pm\sqrt{b^2-4ac}}{2a}.
\tag{1.1.2}
\end{equation}
当上式右边的分子不为零时，也可以使用平方差公式，把根式变到分母上：
\begin{equation}
x_{\pm}=\frac{2c}{-b\mp\sqrt{b^2-4ac}}.
\tag{1.1.3}
\end{equation}
从解析上讲，(1.1.2) 和 (1.1.3) 式是完全等价的，但对于计算机来说就不一定了。取 $b=-1000.001,\ a=c=1$，容易知道，这个方程的精确解为 $x_+=1000,\ x_-=0.001$。如果按照 (1.1.2) 式，通过计算机编程以单精度浮点数做运算，则有（注意：由于不少编程语言存在形如 \texttt{x+ = y} 的语法，我们在下面的算法中以 \texttt{x1}, \texttt{x2} 分别指代 $x_+,x_-$）

\begin{lstlisting}[style=bellacode]
a = 1
b = -1000.001
c = 1
x1 = (-b+sqrt(b^2-4*a*c))/2/a
x2 = (-b-sqrt(b^2-4*a*c))/2/a
write x1
write x2
\end{lstlisting}
运行给出
\begin{lstlisting}[style=bellacode]
1000.00000
1.00708008E-03
\end{lstlisting}
可以发现，$x_+$ 十分精确，而 $x_-$ 与精确结果的误差达到了 $0.7\%$，这在数值计算上是不可忽视的。然而，如果按照 (1.1.3) 式计算，$x_-$ 能比较精确地计算出来，$x_+$ 却又变得不准确了。想要理解我们所观察到的这个现象，我们必须考察计算机对浮点数的表示方法。在计算机中，是通过一串二进制数来近似表示一个实数的。如图\ref{fig:ieee754}所示，在 IEEE-754 规范下，第一个二进制数 $s$ 代表浮点数的正负号，随后一串 $e_1e_2\cdots e_m$ 代表指数，而最后一段 $d_1d_2\cdots d_n$ 代表一个二进制小数，即
\begin{equation}
x=(-1)^s\times 2^{(e_1e_2\cdots e_m)_2-2^{m-1}+1}\times (1.d_1d_2\cdots d_n)_2.
\tag{1.1.4}
\end{equation}

\begin{figure}[htbp]
\centering
\begin{tikzpicture}[x=4.1mm,y=6mm,font=\scriptsize]
\foreach \i in {0,...,31} {\draw (\i,0) rectangle ++(1,1);}
\node at (0.5,0.5) {1};
\foreach \i/\v in {1/1,2/0,3/0,4/0,5/0,6/0,7/1,8/0} {\node at (\i+0.5,0.5) {\v};}
\foreach \i/\v in {9/1,10/0,11/0,12/1,13/1,14/0,15/0,16/0,17/0,18/0,19/0,20/0,21/0,22/0,23/0,24/0,25/0,26/0,27/0,28/0,29/0,30/0,31/0} {\node at (\i+0.5,0.5) {\v};}
\node[above] at (0.5,1.4) {符号位};
\draw[decorate,decoration={brace,mirror,amplitude=3pt}] (1,-0.1)--(9,-0.1) node[midway,below=3pt]{指数位（8位）};
\draw[decorate,decoration={brace,amplitude=3pt}] (9,1.1)--(32,1.1) node[midway,above=3pt]{小数位（23位）};
\end{tikzpicture}
\caption{在计算机里对 $-12.75$ 在单精度浮点系统中的表示}
\label{fig:ieee754}
\end{figure}

上面的表达式相当于用二进制的科学记数法来表示实数。在这个表示法中，指数（对于正规数）被限制在 $2-2^{m-1}$ 至 $2^{m-1}-1$ 的范围内，而小数精确到小数点后 $n$ 位。\jiaokan{原文将正规数的指数范围写为 $-2^{m-1}+1$ 至 $2^{m-1}$。IEEE-754 中全零和全一的指数域分别用于次正规数/零与无穷大/NaN，因此正规数的无偏指数范围应为 $2-2^{m-1}$ 至 $2^{m-1}-1$；单精度为 $-126$ 至 $127$，双精度为 $-1022$ 至 $1023$。式 (1.1.4) 也只适用于正规有限数。}因此，在使用浮点数来存储一个表示范围内的实数时，存在着不大于 $\varepsilon_{\mathrm{mach}}=2^{-n}$ 的相对误差，我们通常将 $\varepsilon_{\mathrm{mach}}$ 称为机器精度。对于单精度浮点数而言，$m=8,n=23$，那么机器精度便约为 $10^{-7}$；对于双精度浮点数而言，$m=11,n=52$，则机器精度约为 $10^{-16}$。因此，在计算机里表示的数，都是用有限字长表示的，除了极少数称为机器数的数能精确表示之外，其余实数都将是近似数。在 IEEE-754 标准中，默认采用最近舍入原则对表示的小数进行舍入\jiaokan{原文此处写“采用最近舍入的原则对表示的小数进行截断”。“round to nearest”是舍入而非截断；默认舍入方式更准确地说是“最近值，遇中点取偶数（roundTiesToEven）”。}，因为计算机有限字长表示所引入的误差，也称为舍入误差。

我们知道，误差在计算过程中将会积累和传播。回到我们刚刚讨论的例子，在
\begin{equation}
-b-\sqrt{b^2-4ac}
\tag{1.1.5}
\end{equation}
这一步中，得到结果的绝对误差约为 $10^{-7}b$，亦即约为 $10^{-4}$。而上述减法的结果本身 $2x_-\sim 10^{-3}$ 却是一个小量。那么我们现在便理解误差是怎么出现的了：用大数减大数得到一个较小的数，得到的绝对误差由大数决定。尽管它相比于大数而言无关紧要，被传播到了这个小数上后，就让结果变得不可靠了。

为了解决上述误差传播所导致的精度损失问题，我们需要避免将两个相近的大数相减，也就是说在求绝对值较小的根时，我们转而使用 (1.1.3) 式。注意到其分母已经在计算绝对值较大的根时算过一次了，我们可以使用 Viète 定理 $x_1x_2=c/a$ 来避免重复计算，从而减少计算量。基于这些考虑，我们可以选择总是先求绝对值较大的根，因此得到如下算法：

\begin{algorithmBox}[算法 1.1\quad 一元二次方程求根]
\begin{lstlisting}[style=bellacode]
function rootQuadratic(a,b,c)
    if(a=0)
        if (b/=0)
            x1 = -c/b
            return x1
        else
            return 'not a proper equation:  a=0, b=0'
        end
    delta = b^2 - 4*a*c
    if(delta>=0)
        if (b/=0)
            x1 = -(b + sgn(b) * sqrt(delta))/2/a
            x2 = c/a/x1
            return x1, x2
        else
            x1 = sqrt(-c/a)
            x2 = - sqrt(-c/a)
            return x1, x2
        end
    else
        return 'no real root'
    end
end
\end{lstlisting}
\end{algorithmBox}

当然，另一个简单易行的提高计算精度的策略是使用双精度来表示浮点数。这样，浮点计算过程中所产生的舍入误差就从 $10^{-7}$ 降低至了 $10^{-16}$。事实上，在 MATLAB 和 Python 这类不强制要求声明变量类型的语言中，浮点数就是默认以双精度的形式进行存储和计算的。但是请读者注意，即便是使用了双精度浮点数，甚至是四精度或者八精度浮点数，不恰当的数值计算过程依旧会导致精度的严重损失，使得最终结果的误差远大于舍入误差。在编写程序的过程中，除了要避免相近的数做减法外，我们还需要避免上溢出、下溢出（过分接近于零的浮点数会损失有效位数）、大数吞小数，以及注意减少运算步骤等等。

\begin{exerciseBox}[练习 1.1\quad 有人宣称调和级数]
\[
\sum_{n=1}^{\infty}\frac{1}{n}
\]
是收敛的，并称在计算机上完成了验证。请根据浮点数的特性说明无法从计算机数值实验得到正确结论的可能原因。
\end{exerciseBox}

\subsubsection{正确与不正确}

前面的例子让我们看到，解析上严格正确的表达式在数值计算中会带来不可忽略的误差，甚至使得数值计算结果完全错误。另一方面，有时解析上不严格的表达式在数值计算中反而可以导向更精确的结果。接下来，我们来讨论这样的一个例子。

我们知道，球 Bessel 函数 $j_l(x)$ 满足三项递归关系
\begin{equation}
j_{l+1}(x)+j_{l-1}(x)=\frac{2l+1}{x}j_l(x),
\tag{1.1.6}
\end{equation}
其中，头两阶的初值由下式解析给出：
\begin{equation}
j_0(x)=\frac{\sin x}{x},\qquad j_1(x)=\frac{\sin x}{x^2}-\frac{\cos x}{x}.
\tag{1.1.7}
\end{equation}
因此，对于任意给定的 $x$，我们原则上可以通过迭代计算出任意 $j_l(x)$ 的数值。那么，我们不妨取 $x=1$，尝试计算 $l=0\sim 15$ 的 $j_l(x)$。所有变量均用单精度浮点数表示，结果如表\ref{tab:bessel-iterations}的“正向迭代”一列所示。

从 $l=7$ 开始，我们就能注意到结果变得不可信了：从解析上我们知道，当 $x/l\ll 1$ 时，$j_l(x)\sim(x/l)^l$，从而 $j_l(x)$ 应该随着 $l$ 增长逐步减小才对，但是数值结果从 $l=7$ 开始是迅速增大的，表现为发散的行为。

那么，我们需要调整正向迭代的算法设计思路，转而探讨反向迭代算法的可能性。反向迭代的难点在于，对于一个较大的 $l$，我们不容易给出其迭代所需要的初值。基于前述解析分析的判断，我们可以直接取 $j_{16}(x)=0$，而将 $j_{15}(x)$ 赋值为任意一个较小的正数。然后我们反向利用迭代公式，将较小 $l$ 的球 Bessel 函数都用 $j_{15}(x)$ 的赋值表示出来。最后，鉴于 $j_0(x)$ 可精确求得，再利用 $j_0(x)/j_{15}(x)$ 的比值，便可确定出 $j_{15}(x)$ 到底该是多少，从而重新归一化定出各个 $j_l(x)$ 的数值。我们将这种办法得到的结果列在表\ref{tab:bessel-iterations}第三列的“反向迭代”中。

显然，反向迭代的初始值是错误的，因为 $j_{16}(x)$ 再怎么小也不会是零。但是，与最后一列的精确值相比，“错误”的方案却给出了非常接近精确值的数值结果。这就非常有趣了：“正确”的不准确，“不正确”的反而准确。这是为什么呢？

\begin{table}[htbp]
\centering
\caption{利用正向迭代和反向迭代算法计算头16个球 Bessel 函数 $j_l(x)$ 在 $x=1$ 时的数值}
\label{tab:bessel-iterations}
\small
\begin{tabular}{rlll}
\toprule
$l$ & 正向迭代 & 反向迭代 & 精确值\\
\midrule
0  & 0.841470957 & 0.841470957 & 0.841471\\
1  & 0.301168680 & 0.301168680 & 0.301169\\
2  & 6.20350838E-02 & 6.20350577E-02 & 6.20351E-2\\
3  & 9.00673866E-03 & 9.00658220E-03 & 9.00658E-3\\
4  & 1.01208687E-03 & 1.01101596E-03 & 1.01102E-3\\
5  & 1.02043152E-04 & 9.25611603E-05 & 9.25612E-5\\
6  & 1.10387802E-04 & 7.15693659E-06 & 7.15694E-6\\
7  & 1.33299828E-03 & 4.79013465E-07 & 4.79013E-7\\
8  & 1.98845863E-02 & 2.82649903E-08 & 2.82650E-8\\
9  & 0.336704969 & 1.49137658E-09 & 1.49138E-9\\
10 & 6.37750959 & 7.11655318E-11 & 7.11655E-11\\
11 & 133.591003 & 3.09955208E-12 & 3.09955E-12\\
12 & 3066.21558 & 1.24166262E-13 & 1.24166E-13\\
13 & 76521.7969 & 4.60463797E-15 & 4.60464E-15\\
14 & 2063022.25 & 1.58957432E-16 & 1.58958E-16\\
15 & 59751124.0 & 5.12765912E-18 & 5.13269E-18\\
\bottomrule
\end{tabular}
\end{table}

事实上，这个问题可以这样理解：从差分方程的相关理论，我们知道递归关系 (1.1.6) 有两个线性无关的解，即球 Bessel 函数 $j_l(x)$ 和球 Neumann 函数 $n_l(x)$。随着 $l$ 的增长，前者将最终逐步趋于零，而后者将发散到无穷大。那么，在迭代过程中，由于舍入误差的存在，我们得到的实际上是两类函数的某个线性组合
\begin{equation}
s_l(x)=\alpha_l j_l(x)+\beta_l n_l(x),
\tag{1.1.8}
\end{equation}
其中，相邻 $\alpha_l,\beta_l$ 之间的差别约为浮点精度。若我们希望通过正向迭代计算某个 $j_l(x)$，那么一个非常小的 $\beta_l$ 就能让结果变得不可信；相反，如果进行反向迭代，并且采取截断令某个 $s_{l_{\max}}(x)=0$，那么随着 $l$ 的减小，因为截断引入的 $\beta_l n_l(x)$ 会越来越小，我们反而能够得到一个准确的结果。实际上，在计算其他类型的递归数列时，我们也会遇见类似的情况，我们务必事先考察递归数列的性质并设计合适的迭代方法。

\begin{exerciseBox}[练习 1.2]
设计正向与反向迭代算法，计算积分
\[
E_n=\int_0^1 x^n e^{x-1}\,dx,\qquad n=0,1,\ldots,20,
\]
对比你所得到的结果。（提示：利用分部积分，我们可以得到正向递推关系式 $E_n=1-nE_{n-1},\ E_0=1-1/e$。）
\end{exerciseBox}

对于一个数值算法，如果输入数据的微小变化不会在输出结果上被过分放大，那么我们称这个算法是稳定的。稳定性包含问题的稳定性和算法的稳定性两个方面。许多混沌问题是天生对初始条件极为敏感的，那么我们不能强求找寻一个稳定的数值算法；而计算一个解析函数的值则通常不存在这样的问题，我们不希望看到结果在随机跳变。

输入数据 $x$ 的扰动对问题解 $y$ 的影响程度的大小称为问题的敏感性，一般可以用\textbf{条件数}来定量描述：
\begin{equation}
\mathrm{cond}\equiv\left|\frac{dy/y}{dx/x}\right|.
\tag{1.1.9}
\end{equation}
条件数表征了解的相对变化与输入数据的相对变化之比。小条件数问题称作不敏感的或者良态的，而大条件数问题称作敏感的或者病态的。

\subsubsection{简单与复杂}

另一些时候，两个在解析上等价的数学表达式，在利用计算机进行数值计算时，都不会对最终结果的精度造成太大的影响，但数值实现所需要的计算量却相差甚远。作为一个经典的例子，我们考虑在给定 $x$ 和 $n$ 时，对下述代数多项式进行数值计算：
\begin{equation}
P_n(x)=\sum_{i=0}^{n}a_i x^i=a_nx^n+a_{n-1}x^{n-1}+\cdots+a_1x+a_0,
\tag{1.1.10}
\end{equation}
其中 $a_i$ 为实常数。若直接逐项计算，第 $i$ 项需要计算 $i$ 次乘法，全加在一起就是 $(n+1)n/2$ 次，另需 $n$ 次加法。稍微聪明一点的做法是，把每一个算出来的 $x^i$ 存下来，用来计算下一步的 $x^{i+1}=x\times x^i$，这样总的乘法次数可以减少到 $2n-1$ 次。实际上，如果我们引入递归序列 $\{p_i\}$：
\begin{equation}
\begin{aligned}
p_n&=a_n,\\
p_i&=xp_{i+1}+a_i,\qquad i=n-1,\ldots,1,0,
\end{aligned}
\tag{1.1.11}
\end{equation}
容易检验，有 $p_0=P_n(x)$。而如此计算所有的 $p_i$ 仅需要 $n$ 次乘法和 $n$ 次加法。这种算法，不仅大大降低了计算量，也必将因为计算次数的减少而降低误差传播。这便是著名的秦九韶算法［西方称为 Horner 算法］。

从解析上看，多项式的显式表达式 (1.1.10) 显然要比迭代形式 (1.1.11) 来得简单直接，但是对于计算机而言，却是迭代形式更加便于计算。这个例子告诉我们，即使对于一个良态的问题，也需要精心设计数值算法，才能更快更准地得到符合精度要求的数值结果。历史上有很多著名的算法，都源于对最优算法设计孜孜不倦的追求，一个大名鼎鼎的例子就是快速 Fourier 变换，将 $O(n^2)$ 的运算量降低为 $O(n\log n)$。

\begin{exerciseBox}[练习 1.3]
设有函数序列 $\{P_k(x)\}$ 满足三项递归关系
\[
P_{k+1}(x)=\alpha_k(x)P_k(x)+\beta_k(x)P_{k-1}(x),
\]
且 $P_0(x),P_1(x)$ 已知，$\{\alpha_k(x),\beta_k(x)\}$ 可以简单计算。试设计算法，高效计算求和式
\begin{equation}
\sum_{k=0}^{n}a_kP_k(x).
\tag{1.1.12}
\end{equation}
\end{exerciseBox}

% R01 stops at the end of the source book's printed page 10.

\subsubsection{小结}

从以上三个简单的例子我们看到，数值分析算法的设计与大家所熟知的解析理论分析有着很大的不同。即使是针对一个简单的数学问题，我们也必须考虑数值算法的方方面面，力求设计出一个好的算法。一个好的数值算法本质上要考虑以下三个方面：\textbf{快、准、稳}。具体来说，快，是指计算所需的时间消耗和存储代价小、计算效率高；准，是指结果精确、误差小；稳，是指算法稳定，对于任何合理输入都能给出准确输出，适用推广。一个数值算法用高级语言编写之后，在正式使用前应该广泛地、多次地进行运行调试，证实算法的稳定性和可靠性。由于程序设计都是模块化的，因此我们应该在处理每一类问题、编写每一个程序时，都养成良好的习惯，才能为数值求解一个大而复杂的问题打下坚实的基础。

\section{第一个程序：开平方}

相信大家都使用过计算器，不少同学可能会有这样的经历：随便输入一个数，再写一个包含 ``Ans''\footnote{这个符号为 ``answer'' 的头三个字母，在计算器中通常表示上一次的输出结果。} 的表达式，随后反复地按 ``='' 键。这样就能看到计算器的输出不断跳动，可能是结果的前几位率先不再变化，然后变化的位数越来越少，最终不再变化；抑或是数值越来越大，直至超出了计算器对数的表示和处理范围，界面上跳出 ``Inf'' 或者 ``NaN''\footnote{``Inf'' 表示无穷大，``NaN'' 表示 ``不是一个数''（not a number）。}；偶尔也会陷入一个奇妙的循环之中，在几个固定的数值之间做周期变化；如果你足够幸运，还会碰见似乎完全没有规律的随机跳动。

这类“游戏”有一个专门的术语，叫作\textbf{迭代}。那些一旦到达就不再变化的点，被称为\textbf{不动点}，而到达不动点的过程就叫作\textbf{迭代收敛}。我们会自然地提出以下问题：
\begin{enumerate}
\item 什么样的迭代会有不动点，有多少个？
\item 所有的不动点在适当的迭代初值下都能收敛吗？
\item 收敛快慢如何？换言之，在计算器上按多少次 ``='' 键可以得到一个不变的结果？
\end{enumerate}

如果在已知的初值范围内，某个迭代只能收敛到一个不动点，并且收敛还相当快，那么我们就可以开心地用它来构造这样一个迭代，让它的不动点正好是某个实际问题的答案，然后从某个初值出发迭代若干次，我们就得到了这个问题的解。数值计算中相当广泛的算法事实上都可以归结于迭代过程。本节以一个著名的开平方的迭代算法为例，来介绍迭代算法中的诸多概念。

\subsection{开平方算法}

开平方的算法有很多，一个比较知名的算法具有如下的迭代格式：
\begin{equation}
 x_{n+1}=\frac12\left(x_n+\frac{a}{x_n}\right).
\tag{1.2.1}
\end{equation}
令 $x^*=x_{n+1}=x_n$，我们很容易就能解出 (1.2.1) 式所对应的不动点 $x^*=\pm\sqrt a$。因此，这个迭代能够给出参数 $a$ 的平方根。为了方便讨论，我们仅考虑 $a>0$ 的情形。

首先我们需要注意一个事实：$x_{n+1}$ 与 $x_n$ 同号。这让我们可以不失一般性地假定 $x_0>0$。进一步，对原式右边进行配方恒等操作，可得到
\begin{equation}
 x_{n+1}=\frac12\left(\sqrt{x_n}-\sqrt{\frac{a}{x_n}}\right)^2+\sqrt a.
\tag{1.2.2}
\end{equation}
上式说明，不论 $x_n$ 是多少，新得到的 $x_{n+1}$ 一定会大于 $\sqrt a$。将最右边的 $\sqrt a$ 移项到左边，进一步变形可得到
\begin{equation}
 \frac{x_{n+1}-\sqrt a}{x_n-\sqrt a}=\frac{x_n-\sqrt a}{2x_n}.
\tag{1.2.3}
\end{equation}
从这个式子可以看出，对于 $x_n>\sqrt a$ 的情形，易证上述比值位于 $\left(0,\frac12\right)$ 间，故 $x_{n+1}$ 依旧大于 $\sqrt a$，且相比 $x_n$ 而言一定距离 $\sqrt a$ 更近。在高等数学中我们学过一个定理：单调有界序列一定存在极限。将它应用到此处，注意到极限值一定是不动点，我们便有结论：对于任意给定的迭代初值 $x_0>0$，序列 $\{x_n\}$ 都会最终收敛到不动点 $\sqrt a$。而对于 $x_0<0$ 的情况，类似地可以证明其会收敛到负根 $-\sqrt a$。像这种不论初值怎么取都能收敛到想要的不动点的情况，我们称之为\textbf{全局收敛}。这样，我们就明白了为什么迭代式 (1.2.1) 可以作为一个开平方算法来使用。

本节一开头提出的问题还差一个没有回答：迭代收敛的快慢如何？记 $x_n$ 与不动点之间的差为 $\varepsilon_n=x_n-x^*$，则 (1.2.3) 式可以改写为 $n\to\infty$ 下的极限形式
\begin{equation}
 \lim_{n\to\infty}\frac{\varepsilon_{n+1}}{\varepsilon_n^2}=\frac{1}{2\sqrt a},
\tag{1.2.4}
\end{equation}
也就是说，当迭代次数足够多之后，每步的误差正比于上一步误差的平方，我们称其为\textbf{二次收敛}。由于迭代充分多次后误差总是远小于 1 的，这是一个相当迅速的收敛速度。

看来，所有的准备工作都完成了，经过考察，我们找到了一个确保收敛、收敛迅速的开平方的算法，我们终于可以开始编写本书的第一个小程序了。但是，还有一点需要考虑：我们得让计算机知道什么时候该停下来，不能陷入无休止的循环。通常来说，我们期待程序给出的结果与真实值 $\sqrt a$ 相差小于某个绝对误差 $\varepsilon$。然而，由于我们并不知道真实值是多少，绝对误差是无法计算的。因此在实践上，常常利用相邻两步迭代值之差的绝对值 $|x_n-x_{n-1}|<\varepsilon$ 来作为对误差的粗略估计和判停标准。\footnote{需要注意的是，对于收敛速度极为缓慢的迭代格式来说，这个误差估计和判停标准显然不尽合理，此时需要极为小心。}以下就是我们的算法。

\begin{algorithmBox}[算法 1.2\quad 开平方]
\begin{lstlisting}[style=bellacode]
function my_sqrt(a,eps)
    if (a<0)
        write 'no real root'
        stop
    endif
    if (a=0)
        return 0
    endif
    x0 = a
    x1 = (x0+a/x0)/2
    while (abs(x1-x0)>eps) do
        x0 = x1
        x1 = (x0+a/x0)/2
    end
    return x1
end
\end{lstlisting}
\end{algorithmBox}
\jiaokan{原算法只判断 $a<0$，随后直接取 $x_0=a$ 并计算 $a/x_0$；当 $a=0$ 时会出现 $0/0$。此处补入 $a=0$ 时直接返回 $0$ 的分支。}

这段伪代码利用一个 while 循环实现了迭代和迭代终止的判断，给出任意正数 $a$ 误差不大于 eps 的平方根。为了方便，我们直接选取 $a$ 本身作为迭代初值。实际上，开平方早已作为基本指令在中央处理器（CPU）上实现了，在多数编程语言中都可以直接通过函数 sqrt() 来调用。CPU 内部所使用的其实也就是这个算法，但会先利用浮点数特殊的存储结构来得到一个相当精确的迭代初值，以减少所需要的迭代次数。\footnote{对这个问题感兴趣的读者可以在互联网上搜索相关内容。}本书中所讨论的数值算法不会深入到这个层次，但针对特定的问题恰当地选择迭代初值这一思路将会在后续章节中反复出现，读者需要时刻牢记初值选取的重要性，初值也是影响计算复杂度的关键因素之一。

下面，利用我们写好的程序来考察两个具体的算例：分别计算 $\sqrt{100}$ 和 $\sqrt2$，将每一步的值以及与标准值之间的相对误差 $\delta_n\equiv x_n/\sqrt a-1$ 输出如表 1.2。\footnote{在没有特殊声明的情况下，本书所有的算例使用的都是双精度浮点数。}可见该算法数步之内就能够得到一个接近机器精度的结果。注意对于 $\sqrt2$ 的计算，在第五步之后相对误差就变成了一个负值且不再变化，与之前解析分析所得的结论不符。这源于前面所提到的重要事实，即计算机是以有限字长的离散浮点数来近似表示连续的实数的。

\begin{table}[htbp]
\centering
\caption{开平方算法的两个算例的迭代序列}
\label{tab:sqrt-iterations}
\small
\begin{tabular}{rrrr}
\toprule
$x_n$ & $\delta_n$ & $x_n$ & $\delta_n$\\
\midrule
100.000 & 9.0 & 2.00000 & 0.41\\
50.5000 & 4.1 & 1.50000 & 0.061\\
26.2401 & 1.6 & 1.41667 & 0.0017\\
15.0255 & 0.50 & 1.41422 & 1.5E-6\\
10.8404 & 0.084 & 1.41421 & 1.1E-12\\
10.0326 & 0.0033 & 1.41421 & $-2.2$E-16\\
10.0001 & 5.3E-6 & 1.41421 & $-2.2$E-16\\
10.0000 & 1.4E-11 & 1.41421 & $-2.2$E-16\\
\bottomrule
\end{tabular}
\end{table}

\subsection{不动点迭代}

一般而言，迭代可以由递归式
\begin{equation}
 x_{n+1}=f(x_n)
\tag{1.2.5}
\end{equation}
来定义。这里的 $x_n$ 可以是整数、实数甚至是向量。这里我们暂时仅考虑 $x_n\in\mathbb R$ 的情形。那么不动点就是那些满足等式 $x^*=f(x^*)$ 的点 $x^*$。我们来考察迭代序列 $x_n$ 向不动点 $x^*$ 收敛的过程。与之前类似，引入 $\varepsilon_n=x_n-x^*$，我们有
\begin{equation}
 \varepsilon_{n+1}=f(x^*+\varepsilon_n)-x^*=f'(x^*)\varepsilon_n+O(\varepsilon_n^2).
\tag{1.2.6}
\end{equation}
对于 $|\varepsilon_n|\ll1$ 且 $f'(x^*)\ne0$ 的情形，我们可以忽略余项 $O(\varepsilon_n^2)$，那么，根据 (1.2.6) 式，每一步的误差大约是上一步的 $f'(x^*)$ 倍。如果 $|f'(x^*)|>1$，则误差是不断增长的，序列会离这个不动点 $x^*$ 越来越远，直至余项不可忽略甚至前述近似完全失效，此时我们称这个不动点是\textbf{不稳定的}，序列无法收敛。对于 $|f'(x^*)|<1$ 的情形，反复迭代可以得到 $\varepsilon_n\approx\varepsilon_0[f'(x^*)]^n$，误差按指数函数形式趋于零。这样的不动点则是\textbf{稳定的}。这样，我们就得到了一个关于迭代的收敛判据，即考察不动点处导数的绝对值的大小。

更一般地，对于任意一个收敛的迭代序列 $\{x_n\}$，若存在常数 $p\ge1$ 和 $C\ge0$，使得极限关系式
\begin{equation}
 \lim_{n\to\infty}\frac{|\varepsilon_{n+1}|}{|\varepsilon_n|^p}=C
\tag{1.2.7}
\end{equation}
成立，则称迭代序列 $\{x_n\}$ 是 $p$ 阶收敛的。其中 $p$ 为该序列的收敛阶，$C$ 为收敛率。

当判断一个迭代过程的收敛快慢时，我们会优先比较它们的收敛阶 $p$。收敛阶越高的算法总是收敛得越快，至少在迭代次数充分多之后是这样。对于 $p=1$ 的情况，我们会尤其关心收敛率 $C$ 的大小。这个时候，我们称 $C\in(0,1)$ 为\textbf{线性收敛}，$C=0$ 为\textbf{超线性收敛}，$C=1$ 为\textbf{次线性收敛}。当 $p=1$ 时，若有 $C>1$，则迭代是无法收敛的。

\begin{exerciseBox}[练习 1.4]
有若干迭代算法，其误差 $\varepsilon_n$ 随迭代次数 $n$ 的变化关系分别为
\[
 e^{-\alpha n},\qquad e^{-n^2/\sigma^2},\qquad \frac{1}{n^2},\qquad 2^{-3^n},
\]
试分别求出各自的收敛率和收敛阶。
\end{exerciseBox}

\subsection{小结}

在本节中，我们从计算器上的“小游戏”出发，引入了数值算法中最重要的思想之一：迭代。以一个开平方算法为例，我们介绍了迭代不动点、收敛性、收敛速度等概念。在本章末尾，读者将面临本书第一个大作业，利用这一节学习的知识以及题目中给出的引导进行数值和物理上的自主探索。事实上，本书的内容都会以这种方式进行编排：每一节从读者熟悉的某个物理或者数学问题出发，引入相关的数值算法，并讨论其一般化的推广；在每一章的末尾，会设置一个课题，读者需要综合运用该章及之前章节中学到的数值计算知识上机编程，并探讨其背后所隐藏的数学与物理规律。

\section*{大作业：逻辑斯谛映射}
\addcontentsline{toc}{section}{大作业：逻辑斯谛映射}

我们来考察一个非常简单而广为人知的数值模型，其由如下迭代关系定义：
\begin{equation}
 x_{n+1}=f(x_n)=r x_n(1-x_n),
\tag{1}
\end{equation}
其中 $r>0$ 为可调参数。(1) 式右端的函数 $f(x)$ 被称为逻辑斯谛（logistic）函数，式中所规定的从 $x_n$ 到 $x_{n+1}$ 的迭代称为逻辑斯谛映射。这一看似简单的二次多项式映射数学模型却描述了物理、生物、化学、工程等众多学科领域中的很多丰富而复杂的非线性动力学行为，感兴趣的同学可以参考相关资料。\footnote{例如：May R M. Simple mathematical models with very complicated dynamics. \textit{Nature}, 1976, 261: 459；Strogatz S H. \textit{Nonlinear Dynamics and Chaos: With Applications to Physics, Biology, Chemistry, and Engineering}. 2nd ed. CRC Press, 2015.}

\begin{enumerate}
\item 作为最初步的认识，分别取 $r=0.5$ 和 $r=1.5$，任取几个 0 到 1 之间的初值 $x_0$，计算序列 $\{x_n\}$。绘图观察并描述它们的行为。
\item 显然，对于不同的 $r$ 值，$x_n$ 将迅速收敛于某一个特定的 $x^*$ 处，而与初值无关。$x^*$ 必将满足如下自治方程：
\begin{equation}
 x^*=f(x^*).
\tag{2}
\end{equation}
作为一个二元二次方程，其存在两个根。试判断并证明哪一个根才是最终收敛的不动点，并画出 $x^*$ 随 $r$ 的变化关系图。一般而言，其收敛阶 $p$ 和收敛率 $C$ 各是多少？
\item 当 $r$ 大于某个特定值 $r_1$ 时，上述条件无法满足。取 $r=r_1+0.1$ 以及不同的初值 $x_0$，计算序列 $\{x_n\}$。绘图观察并描述它们的行为。
\item 序列终将在某两个 $x_1^*$ 和 $x_2^*$ 之间来回振荡。事实上，考察复合迭代 $x_{n+2}=f(f(x_n))$，其所定义的序列依旧将收敛于某一个固定值，从而依旧可以使用前述方法来分析。试证明，这类迭代收敛的必要条件是 $|f'(x_1^*)f'(x_2^*)|\le1$，并在第 2 问的图中补上 $x_1^*$ 和 $x_2^*$ 随 $r$ 的变化关系。
\item 继续缓慢增大 $r$，请展示周期为二的振荡会逐渐变成周期四、周期八、……。试定义一个量，可以一般性地描述序列向各类稳定振荡的平均收敛速度，并可由一般有限序列近似算出。在 $r\in(0,4)$ 上等距取点，绘制收敛速度随 $r$ 的变化，描述其与不同振荡周期的关系。
\item 在第 4 问的图中补上后续振荡值随 $r$ 的变化关系。依次缩小你的作图范围，将坐标原点分别放在周期一到周期二的分岔点、周期二到周期四的分岔点、周期四到周期八的分岔点，等等。描述你所观察到的现象。
\item 计算相邻分岔点之间的横轴距离 $\Delta r$，说明相邻 $\Delta r$ 之比渐近于一常数 $F$，并给出 $F$ 以及无穷周期分岔点 $r_\infty$ 的值。
\item 对于 $r=4$ 的情形，试解析求解该迭代序列，论证此时一般不存在稳定的振荡周期（提示：做代换 $x=\sin^2 y$）。
\item 选择另一个你喜欢的函数 $g(x)$，其在从原点开始的某段区域上是凸函数，将 $x(1-x)$ 替换成 $g(x)$，重复上面第 1 至 7 问的计算，说明你能看到类似的现象，并得到完全相同的 $F$ 值。
\end{enumerate}

% R02 extends through the end of Chapter 1 (source printed page 16).
